\documentclass[11pt,letterpaper]{article}

% ── geometry & layout ────────────────────────────────────────────────
\usepackage[margin=1in]{geometry}
\usepackage{fancyhdr}
\usepackage{setspace}
\onehalfspacing

% ── math ─────────────────────────────────────────────────────────────
\usepackage{amsmath,amssymb,amsthm,mathtools}
\usepackage[T1]{fontenc}

% ── colour & boxes ───────────────────────────────────────────────────
\usepackage[dvipsnames]{xcolor}
\usepackage[most]{tcolorbox}
\usepackage{enumitem}
\usepackage{booktabs,array,tabularx}
\usepackage{tikz}
\usetikzlibrary{arrows.meta,positioning,calc,decorations.pathreplacing,matrix}
\usepackage[colorlinks=true,linkcolor=MidnightBlue,citecolor=OliveGreen,urlcolor=BrickRed]{hyperref}

% ── theorem environments ─────────────────────────────────────────────
\theoremstyle{definition}
\newtheorem{definition}{Definition}[section]
\newtheorem{example}{Example}[section]
\theoremstyle{plain}
\newtheorem{theorem}{Theorem}[section]
\newtheorem{proposition}{Proposition}[section]
\newtheorem{lemma}{Lemma}[section]
\newtheorem{corollary}{Corollary}[section]
\theoremstyle{remark}
\newtheorem*{remark}{Remark}

% ── macros (match paper/preamble.tex) ────────────────────────────────
\newcommand{\R}{\mathbb{R}}
\newcommand{\E}{\mathbb{E}}
\newcommand{\D}{\mathcal{D}}
\newcommand{\Pset}{\mathcal{P}}
\newcommand{\agents}{N}
\newcommand{\votes}{v}
\newcommand{\budget}{B}
\newcommand{\cost}{c}
\newcommand{\type}{\theta}
\newcommand{\typespace}{\Theta}
\newcommand{\outcomes}{X}
\newcommand{\scf}{f}
\newcommand{\DEM}{\text{DEM}}
\newcommand{\UAM}{\text{UAM}}
\newcommand{\MES}{\text{MES}}
\newcommand{\mesrho}{\varrho}      % MES per-supporter cap; never \rho
\newcommand{\salience}{\alpha}
\DeclareMathOperator*{\argmin}{arg\,min}
\DeclareMathOperator*{\argmax}{arg\,max}

% ── citation fallbacks (standalone, no .bib) ─────────────────────────
\providecommand{\citet}[1]{Peters, Pierczy\'nski \& Skowron (2021)}
\providecommand{\citep}[1]{(Peters, Pierczy\'nski \& Skowron 2021)}
\providecommand{\citealp}[1]{Peters, Pierczy\'nski \& Skowron 2021}

% ── pedagogical boxes ────────────────────────────────────────────────
\newtcolorbox{keyidea}{colback=OliveGreen!6,colframe=OliveGreen!65,
  fonttitle=\bfseries,title={Key idea},boxsep=2pt,left=5pt,right=5pt,top=3pt,bottom=3pt}
\newtcolorbox{intuition}{colback=MidnightBlue!5,colframe=MidnightBlue!60,
  fonttitle=\bfseries,title={Intuition},boxsep=2pt,left=5pt,right=5pt,top=3pt,bottom=3pt}
\newtcolorbox{connection}{colback=TealBlue!6,colframe=TealBlue!60,
  fonttitle=\bfseries,title={Connection to Durango},boxsep=2pt,left=5pt,right=5pt,top=3pt,bottom=3pt}
\newtcolorbox{naming}{colback=Violet!6,colframe=Violet!65,
  fonttitle=\bfseries,title={A name we are introducing (our label, not standard terminology)},
  boxsep=2pt,left=5pt,right=5pt,top=3pt,bottom=3pt}
\newtcolorbox{watchout}{colback=BrickRed!5,colframe=BrickRed!60,
  fonttitle=\bfseries,title={Watch out --- a subtle point that is easy to get wrong},
  boxsep=2pt,left=5pt,right=5pt,top=3pt,bottom=3pt}
\newtcolorbox{aside}{colback=Dandelion!8,colframe=Dandelion!80!black,
  boxsep=2pt,left=5pt,right=5pt,top=3pt,bottom=3pt}

% ── header ───────────────────────────────────────────────────────────
\pagestyle{fancy}\fancyhf{}
\fancyhead[L]{\small\textit{Working Notebook, v2 --- combined \& corrected}}
\fancyhead[R]{\small\textit{Mechanism Design for Participatory Budgeting}}
\fancyfoot[C]{\thepage}
\renewcommand{\headrulewidth}{0.4pt}

\title{%
  \vspace{-0.8cm}
  {\large\scshape Working Notebook --- Second Edition (combined \& corrected)}\\[0.6em]
  {\LARGE\bfseries Mechanism Design for\\[0.2em] Participatory Budgeting}\\[0.5em]
  {\large From first principles to the Method of Equal Shares,\\[0.2em]
  with full motivations and worked examples}
}
\author{\textbf{Mario Alberto Garc\'ia Meza}\\[2pt]
  \small Facultad de Econom\'ia, Contadur\'ia y Administraci\'on\\
  \small Universidad Ju\'arez del Estado de Durango (FECA--UJED)}
\date{June 2026 \;\textbullet\; Self-study \& writing companion}

\begin{document}
\maketitle
\thispagestyle{fancy}

\begin{abstract}
\noindent This notebook rebuilds, in a single volume, the theory behind our paper on
Durango's participatory budgeting. It starts from the basics of mechanism design, develops
the formal model of the city's voting rule, proves the propositions that diagnose what goes
wrong, introduces the Method of Equal Shares as a remedy, and closes with how we will test
everything in simulation. It is written for an economist who knows this material but has been
away from it for a while: every result comes with motivation, intuition, and a worked example.
It supersedes the three earlier working notebooks, folding in every correction from the
mathematical review.
\end{abstract}

\begin{aside}
\small\textbf{How to read this, and three conventions.}
\textbf{(1) Boxes.} Green boxes hold the load-bearing idea; blue boxes give intuition; teal
boxes connect the abstract result to Durango; \emph{red ``Watch out'' boxes} flag a subtle
point that is genuinely easy to get wrong (these are where the corrections from the review
live). \textbf{(2) Names.} Some terms are standard and cited; others are labels we are coining
for this project, and these are always introduced explicitly in a violet ``A name we are
introducing'' box, so you can always tell ours from the literature's. In particular, the
\emph{Method of Equal Shares} and \emph{Extended Justified Representation} are established
(\citealp{peters2021}); the \emph{District-Exclusivity Mechanism} and \emph{Unrestricted
Allocation Mechanism} are our labels. \textbf{(3) Notation.} We follow \texttt{paper/preamble.tex};
the only trap is that the Method of Equal Shares' per-supporter cost cap is written $\mesrho$
(\texttt{\textbackslash mesrho}), never $\rho$, which clashes with a built-in.
\end{aside}

\tableofcontents
\newpage

% ═════════════════════════════════════════════════════════════════════
\section{What mechanism design is, and why we need it}
\label{sec:what}
% ═════════════════════════════════════════════════════════════════════

Start with the problem a city government actually faces. Durango wants to spend a fixed pot of
public money on the neighbourhood projects its citizens care about most. The trouble is that
the government does not know how much people care. That information --- who values a new park,
who needs the drainage fixed, how badly --- lives inside people's heads. The only way to get
at it is to \emph{ask}, through some voting or reporting procedure. And the moment you ask,
people may answer strategically: they may exaggerate, coordinate, or game whatever rule you
set up.

\begin{keyidea}
Game theory and mechanism design ask mirror-image questions.
Game theory: \emph{the rules are fixed --- what will people do?}
Mechanism design: \emph{the goal is fixed --- what rules will make people's behaviour produce it?}
We are squarely in the second business. The voting rule \emph{is} the mechanism, and our
question is whether Durango's rule produces good outcomes, and if not, what rule would.
\end{keyidea}

The field was launched by Hurwicz and developed by Maskin and Myerson (the 2007 Nobel). We
will need only its first principles, but we will need them precisely, because the whole
argument of the paper is that Durango's rule fails a specific, nameable list of those
principles.

\begin{connection}
Throughout, keep one concrete picture in mind. You go to vote in the \emph{presupuesto
participativo}. You care about one project --- say, the university's --- and you are handed six
votes and a ballot listing many projects you have never heard of, all of them in a single
district you must commit to up front. What you would \emph{like} to express (``fund this one
thing I know about, across town from where I live'') is not something the ballot lets you say.
That mismatch between what you know and what the rule lets you report is the seed of everything
that follows.
\end{connection}

% ─────────────────────────────────────────────────────────────────────
\subsection{The environment: agents, types, outcomes}
% ─────────────────────────────────────────────────────────────────────

We build the machinery one piece at a time. Every definition here is used later, so it is
worth going slowly.

\begin{definition}[Environment]
\label{def:environment}
An \textbf{environment} consists of:
\begin{enumerate}[label=(\roman*),nosep]
  \item a finite set of \textbf{agents} $\agents=\{1,\dots,n\}$;
  \item for each agent $i$, a set of possible \textbf{types} $\typespace_i$; a type
        $\type_i\in\typespace_i$ encodes $i$'s private information. The set of type profiles is
        $\typespace=\typespace_1\times\cdots\times\typespace_n$;
  \item a set of \textbf{outcomes} $\outcomes$;
  \item for each agent $i$, a \textbf{utility function} $u_i:\outcomes\times\typespace_i\to\R$.
\end{enumerate}
We write $\type_{-i}=(\type_1,\dots,\type_{i-1},\type_{i+1},\dots,\type_n)$ for everyone's type
but $i$'s, so a profile is $\type=(\type_i,\type_{-i})$.
\end{definition}

\begin{intuition}
The \emph{type} is the whole source of difficulty. It is what the agent knows and the planner
does not. In an auction your type is your valuation; in voting it is your preferences. If types
were public there would be no problem --- the planner would just compute the best outcome. Every
hard question in this notebook traces back to types being private.
\end{intuition}

\begin{connection}
For us the agents are Durango's voters, the outcomes are sets of funded projects (within
budget), and the utility is how much a voter benefits from a given funded set. But our type
will carry an extra dimension that standard models omit: not just \emph{how much} a voter values
each project, but \emph{whether they know enough to value it at all}. We return to this in
Section~\ref{sec:model}; flag it now as the thing that makes our model richer than the textbook.
\end{connection}

% ─────────────────────────────────────────────────────────────────────
\subsection{Social choice functions: what the planner would do if she knew everything}
% ─────────────────────────────────────────────────────────────────────

\begin{definition}[Social choice function]
\label{def:scf}
A \textbf{social choice function} (SCF) is a map $\scf:\typespace\to\outcomes$. It says: if the
planner knew the true type profile $\type$, she would choose outcome $\scf(\type)$.
\end{definition}

\noindent Two SCFs are worth naming because they bracket the equity debate that runs through
the paper.

\begin{example}[Utilitarian SCF]
\label{ex:util}
$\scf^{\text{util}}(\type)=\argmax_{x\in\outcomes}\sum_{i}u_i(x,\type_i)$: pick the outcome
maximising total utility. This is the natural ``fund what people value most in total'' benchmark.
\end{example}

\begin{example}[Rawlsian SCF]
\label{ex:rawls}
$\scf^{\text{Rawls}}(\type)=\argmax_{x\in\outcomes}\min_i u_i(x,\type_i)$: maximise the
worst-off agent's utility. This speaks to the worry about marginalised districts --- even if
total welfare is maximised by funding the city centre, a maximin planner would refuse to leave
a peripheral neighbourhood with nothing.
\end{example}

\begin{connection}
The paper does not have to choose one of these. Its point is subtler: the \emph{current rule}
sacrifices utilitarian efficiency, while a naive fix sacrifices the Rawlsian/equity concern,
and a well-designed rule can respect both. Holding these two SCFs in mind keeps that tension in
view.
\end{connection}

% ─────────────────────────────────────────────────────────────────────
\subsection{Mechanisms: the rules of the game}
% ─────────────────────────────────────────────────────────────────────

\begin{definition}[Mechanism]
\label{def:mechanism}
A \textbf{mechanism} is a pair $(\mathcal{M},g)$ where $\mathcal{M}=M_1\times\cdots\times M_n$
is the \textbf{message space} (agent $i$ sends $m_i\in M_i$) and $g:\mathcal{M}\to\outcomes$ is
the \textbf{outcome function}. A mechanism is \textbf{direct} if $M_i=\typespace_i$ --- agents
simply report a type.
\end{definition}

A mechanism turns the environment into a game: everyone sends a message, and $g$ converts the
messages into an outcome.

\begin{example}[Durango's rule as a mechanism]
\label{ex:durango}
The message of voter $i$ is a pair: a district $d_i$ (one must be chosen) and an allocation of
$\votes$ votes among the projects of that district. The outcome function adds up the votes each
project receives, ranks projects by total votes, and funds them top-down until the budget runs
out. This is \emph{not} a direct mechanism: a voter never reports their preferences over all
projects, only a district and a within-district allocation.
\end{example}

\begin{watchout}
The compression in Example~\ref{ex:durango} is the root of the entire paper. The message space
$M_i$ is dramatically smaller than the type space $\typespace_i$: a voter who cares about
projects in three districts can speak about only one. \textbf{Information is lost at the moment
of reporting, before any votes are even counted.} Most analyses of voting focus on how votes
are aggregated; our focus is one step earlier, on what the ballot does not let you say.
\end{watchout}

% ─────────────────────────────────────────────────────────────────────
\subsection{Dominant strategies and implementation}
% ─────────────────────────────────────────────────────────────────────

When does a mechanism actually deliver the SCF the planner wanted? The strongest answer:

\begin{definition}[Dominant strategy]
\label{def:dominant}
A strategy $s_i^*:\typespace_i\to M_i$ is \textbf{dominant} for $i$ if, for every type $\type_i$
and every message profile $m_{-i}$ of the others,
\[
  u_i\bigl(g(s_i^*(\type_i),m_{-i}),\type_i\bigr)\ \ge\ u_i\bigl(g(m_i,m_{-i}),\type_i\bigr)
  \qquad\text{for all } m_i\in M_i.
\]
That is, $s_i^*(\type_i)$ is optimal no matter what anyone else does.
\end{definition}

\begin{definition}[Implementation in dominant strategies]
\label{def:implementation}
A mechanism $(\mathcal{M},g)$ \textbf{implements} the SCF $\scf$ in dominant strategies if there
is a dominant-strategy profile $s^*=(s_1^*,\dots,s_n^*)$ with $g(s^*(\type))=\scf(\type)$ for all
$\type$.
\end{definition}

\begin{intuition}
Dominant-strategy implementation is the gold standard because it asks nothing of the agents'
beliefs: whatever you think others will do, your best move is $s_i^*$, and when everyone plays
that way the planner gets exactly what she wanted. No coordination, no second-guessing.
\end{intuition}

\begin{connection}
Durango's rule already fails this badly, and the reason is instructive. Your best district
depends on what you expect others to do --- join the big bloc forming in district $A$, or stay in
your home district $B$ where your marginal vote counts for more? A rule whose best play depends
on your forecast of everyone else has no dominant strategies. We will make this precise as a
\emph{coordination game} in Section~\ref{sec:distortions}.
\end{connection}

% ─────────────────────────────────────────────────────────────────────
\subsection{Incentive compatibility}
% ─────────────────────────────────────────────────────────────────────

The honesty version of the same idea:

\begin{definition}[DSIC and BIC]
\label{def:ic}
An SCF $\scf$ is \textbf{dominant-strategy incentive compatible (DSIC)}, or
\textbf{strategy-proof}, if for all $i$, all true $\type_i$, all misreports $\hat\type_i$, and
all $\type_{-i}$,
\[
  u_i\bigl(\scf(\type_i,\type_{-i}),\type_i\bigr)\ \ge\ u_i\bigl(\scf(\hat\type_i,\type_{-i}),\type_i\bigr).
\]
It is \textbf{Bayesian incentive compatible (BIC)} if truth-telling is optimal only in
expectation over $\type_{-i}$ (given a common prior).
\end{definition}

\begin{remark}
DSIC is stronger than BIC: ``best no matter what'' implies ``best on average.'' DSIC is what we
would love to have, because it frees voters from having to model each other. As the next two
sections show, it is also usually impossible --- which is why the paper does not pretend to
deliver it, and is honest about that.
\end{remark}

% ─────────────────────────────────────────────────────────────────────
\subsection{The Revelation Principle}
% ─────────────────────────────────────────────────────────────────────

Here is the result that makes the whole subject tractable: it lets us stop searching over the
infinitely many possible message spaces.

\begin{theorem}[Revelation Principle; Myerson 1981]
\label{thm:revelation}
If some mechanism implements an SCF $\scf$ in dominant strategies, then the \emph{direct}
mechanism in which agents report types and truth-telling is dominant also implements $\scf$.
\end{theorem}

\begin{proof}
Let $(\mathcal{M},g)$ implement $\scf$ with dominant profile $s^*$, so $g(s^*(\type))=\scf(\type)$
for all $\type$. Build a direct mechanism $g':\typespace\to\outcomes$ by
\[
  g'(\hat\type)=g\bigl(s^*(\hat\type)\bigr)=g\bigl(s_1^*(\hat\type_1),\dots,s_n^*(\hat\type_n)\bigr):
\]
it reads reported types, runs them through the equilibrium strategies, and feeds the result to
the original $g$.

\smallskip\noindent\emph{Claim: truth-telling is dominant in $(\typespace,g')$.}
Fix agent $i$ with true type $\type_i$ and any reports $\hat\type_{-i}$ of the others, and set
$m_{-i}'=s_{-i}^*(\hat\type_{-i})$. Telling the truth yields
$u_i\bigl(g(s_i^*(\type_i),m_{-i}'),\type_i\bigr)$, while reporting $\hat\type_i$ yields
$u_i\bigl(g(s_i^*(\hat\type_i),m_{-i}'),\type_i\bigr)$. Because $s_i^*$ is dominant in the
original mechanism, the first is at least the second (apply Definition~\ref{def:dominant} with
$m_i=s_i^*(\hat\type_i)$). So no misreport helps. Finally, when everyone is truthful,
$g'(\type)=g(s^*(\type))=\scf(\type)$.
\end{proof}

\begin{keyidea}
The trick is \emph{simulation}: the direct mechanism quietly plays the old equilibrium strategy
on your behalf. If you could not beat $s_i^*$ by sending a different message, you cannot beat
truth-telling by reporting a different type. The pay-off: to ask ``is this goal achievable?'' we
only ever have to check whether honesty is incentive-compatible in a direct mechanism --- we do
not have to invent clever message spaces.
\end{keyidea}

% ─────────────────────────────────────────────────────────────────────
\subsection{Gibbard--Satterthwaite, and why it is only a backdrop here}
% ─────────────────────────────────────────────────────────────────────

The Revelation Principle says ``check honesty in direct mechanisms.'' The next theorem says that
this check almost always fails.

\begin{theorem}[Gibbard 1973; Satterthwaite 1975]
\label{thm:gs}
Let $|\outcomes|\ge3$. If $\scf:\typespace\to\outcomes$ is surjective and DSIC over an
\emph{unrestricted} domain of preferences, then $\scf$ is \textbf{dictatorial}: there is an agent
$j$ whose top choice is always selected.
\end{theorem}

The full proof is long and combinatorial (it is in Mas-Colell, Whinston \& Green, ch.~23, which we
cite rather than reproduce). The engine of it is one lemma, which is worth proving because it is
clean and you will recognise it elsewhere.

\begin{lemma}[DSIC $\Rightarrow$ monotonicity]
\label{lem:mono}
Say $\scf$ is \textbf{monotone} if, whenever $\scf(\type)=x$ and we change one agent's type from
$\type_i$ to $\type_i'$ in a way that does not push $x$ down in $i$'s ranking (anything below $x$
stays below $x$), then $\scf(\type_i',\type_{-i})=x$. If $\scf$ is DSIC, it is monotone.
\end{lemma}

\begin{proof}
Let $y=\scf(\type_i',\type_{-i})$. DSIC at the true type $\type_i$ gives
$u_i(x,\type_i)\ge u_i(y,\type_i)$, so $y$ sits weakly below $x$ under $\type_i$. Because $x$ did
not fall when we moved to $\type_i'$, $y$ stays weakly below $x$ there too:
$u_i(x,\type_i')\ge u_i(y,\type_i')$. But DSIC at the true type $\type_i'$ gives the reverse,
$u_i(y,\type_i')\ge u_i(x,\type_i')$. Hence $u_i(x,\type_i')=u_i(y,\type_i')$, and with strict
preferences this forces $y=x$.
\end{proof}

\noindent The rest of the argument shows that monotonicity plus surjectivity, pushed across all
pairs of outcomes and all agents, makes some single agent pivotal for everything --- a dictator.

\begin{watchout}
It is tempting to wield Gibbard--Satterthwaite as a club: ``no good voting rule is
strategy-proof, end of story.'' But read the hypothesis: it needs an \emph{unrestricted} domain
over $\ge 3$ outcomes. Our participatory-budgeting domain is \emph{restricted} --- utilities are
additive across projects, outcomes must fit a budget, and projects sit in a district structure.
On restricted domains the theorem simply does not apply, and that is exactly the crack through
which proportional rules like the Method of Equal Shares climb. So in this notebook GS is
\emph{motivation only}: it explains why we should not expect, or promise, strategy-proofness. The
sharp impossibility we actually invoke later (for the Method of Equal Shares specifically) is a
different, domain-appropriate result of \citet{peters2021}. \textbf{Never say ``GS implies the
Method of Equal Shares is manipulable''} --- that conflates two different theorems.
\end{watchout}

\begin{connection}
The lesson for the paper's framing: do not over-promise. The current rule is not bad because it
fails some unreachable ideal of strategy-proofness; \emph{nothing} reachable is
strategy-proof here. It is bad for concrete, fixable reasons --- truncation, noise, coordination,
capture --- which we now make precise, and which a better rule can address even though it, too,
will not be strategy-proof.
\end{connection}

% ═════════════════════════════════════════════════════════════════════
\section{The participatory budgeting model}
\label{sec:model}
% ═════════════════════════════════════════════════════════════════════

Now we specialise the general machinery to Durango. Everything in Section~\ref{sec:what} was
textbook; from here on the objects are ours, built to capture the two features that make this
problem distinctive: voters who only know about part of the ballot, and a city geography that
the rule treats as a hard partition.

% ─────────────────────────────────────────────────────────────────────
\subsection{The environment}
% ─────────────────────────────────────────────────────────────────────

\begin{definition}[Participatory budgeting environment]
\label{def:pb}
A \textbf{participatory budgeting environment} is a tuple
$\langle\agents,\D,\Pset,\budget,\votes,u,\sigma\rangle$ where:
\begin{enumerate}[label=(\roman*),nosep]
  \item $\agents=\{1,\dots,n\}$ is the set of \textbf{voters};
  \item $\D=\{1,\dots,m\}$ is the set of \textbf{districts};
  \item $\Pset=\bigcup_{d\in\D}\Pset^d$ is the set of \textbf{projects}, where $\Pset^d$ are the
        projects of district $d$ and $\Pset^d\cap\Pset^{d'}=\emptyset$ for $d\ne d'$ (each project
        belongs to exactly one district); each project $p$ has a cost $\cost(p)>0$;
  \item $\budget>0$ is the total \textbf{budget};
  \item $\votes\in\mathbb{N}$ is the \textbf{vote endowment} per voter (Durango uses $\votes=6$);
  \item $u_i:\Pset\to\R_{\ge0}$ is voter $i$'s \textbf{utility} over individual projects, extended
        to funded sets by \textbf{additive separability}: $U_i(S)=\sum_{p\in S}u_i(p)$;
  \item $\sigma_i:\Pset\to\{0,1\}$ is voter $i$'s \textbf{information function}: $\sigma_i(p)=1$
        iff $i$ knows enough about $p$ to form a genuine preference, and $0$ otherwise.
\end{enumerate}
The full type of voter $i$ is the pair $\type_i=(u_i,\sigma_i)$ --- preferences \emph{and} information.
\end{definition}

\begin{keyidea}
Two modelling choices do all the work later, so notice them now.
\textbf{(1) The information function $\sigma_i$.} Standard PB models silently assume
$\sigma_i\equiv1$: everyone can evaluate everything. Real voters cannot. Encoding ``do you even
know this project?'' separately from ``how much would you value it?'' is what lets us talk about
\emph{noise} (Prop.~\ref{prop:noise}) and \emph{salience} (Prop.~\ref{prop:equity}). This connects
to Downs' rational ignorance and to costly information acquisition in voting (Martinelli 2006).
\textbf{(2) The disjoint districts.} Because $\Pset^d\cap\Pset^{d'}=\emptyset$, a voter who
``enters'' district $d$ can touch \emph{no} project elsewhere. That single partition is the
structural fact from which all five distortions flow.
\end{keyidea}

\begin{remark}[On additive separability]
$U_i(S)=\sum_{p\in S}u_i(p)$ assumes no complementarity or substitution across projects (a new
park and a new road are valued independently). It is a genuine restriction, but it is standard in
the PB literature and it is what makes welfare bookkeeping tractable. State it honestly as an
assumption in the paper rather than hiding it.
\end{remark}

% ─────────────────────────────────────────────────────────────────────
\subsection{The current rule, named}
% ─────────────────────────────────────────────────────────────────────

\begin{naming}
We need a short name for Durango's actual rule. The literature has no standard term for it, so we
coin one: the \textbf{District-Exclusivity Mechanism}, abbreviated \textbf{DEM}. ``Exclusivity''
because a voter is locked into a single district; nothing about the name is borrowed from
elsewhere --- whenever you see ``DEM'' it means precisely the rule in Definition~\ref{def:dem}.
\end{naming}

\begin{definition}[District-Exclusivity Mechanism, DEM]
\label{def:dem}
Under the \textbf{DEM} each voter $i$ submits a message in two parts: a district $d_i\in\D$, and a
vote allocation $\mathbf{x}_i=(x_{i,p})_{p\in\Pset^{d_i}}$ with $x_{i,p}\ge0$ and
$\sum_{p\in\Pset^{d_i}}x_{i,p}=\votes$. The total vote of project $p$ is
$V(p)=\sum_{i:\,d_i=d(p)}x_{i,p}$, where $d(p)$ is $p$'s district. Projects are ranked by $V(p)$
and funded greedily --- highest first --- until the budget $\budget$ cannot afford the next one.
\end{definition}

\begin{remark}
The DEM is not a direct mechanism (it does not elicit the full type), and its message space
$M_i^{\DEM}=\{(d,\mathbf{x}):d\in\D,\ \mathbf{x}\text{ an allocation of }\votes\text{ over }\Pset^d\}$
is far smaller than the type space. This is the formal restatement of the ``compression'' warning
from Example~\ref{ex:durango}.
\end{remark}

% ─────────────────────────────────────────────────────────────────────
\subsection{The obvious alternative, named}
% ─────────────────────────────────────────────────────────────────────

The first thing anyone proposes on seeing the DEM is ``just drop the district lock-in.'' We will
need that strawman as a benchmark, so we name it too.

\begin{naming}
We call the rule that keeps everything about the DEM \emph{except} the district restriction the
\textbf{Unrestricted Allocation Mechanism}, abbreviated \textbf{UAM}. Again this is our label, not
a term from the literature. Its role in the paper is to be the tempting-but-flawed benchmark: it
fixes the efficiency problems and then fails on equity (Prop.~\ref{prop:equity}).
\end{naming}

\begin{definition}[Unrestricted Allocation Mechanism, UAM]
\label{def:uam}
The \textbf{UAM} is identical to the DEM except that there is no district-selection step: voter
$i$ allocates $\votes$ votes over \emph{all} projects, $\mathbf{x}_i=(x_{i,p})_{p\in\Pset}$ with
$x_{i,p}\ge0$ and $\sum_{p\in\Pset}x_{i,p}=\votes$. The greedy funding rule is unchanged.
\end{definition}

\begin{remark}
Every feasible DEM allocation is also a feasible UAM allocation (the special case where all of a
voter's votes happen to land in one district). So the DEM's reachable vote profiles are contained
in the UAM's --- a fact we use directly in Proposition~\ref{prop:trunc}.
\end{remark}

% ─────────────────────────────────────────────────────────────────────
\subsection{How a voter votes: sincere and informed-sincere}
% ─────────────────────────────────────────────────────────────────────

To compare rules cleanly, we sometimes hold behaviour fixed at a simple, honest benchmark and ask
what the \emph{rule} does. That benchmark is sincere voting.

\begin{definition}[Sincere and informed-sincere voting]
\label{def:sincere}
Under \textbf{sincere voting} a voter spreads votes in proportion to utility over the projects the
rule lets them touch:
\[
  x_{i,p}^{\text{sincere}}=\votes\cdot\frac{u_i(p)}{\sum_{q}u_i(q)},
\]
the sum running over $\Pset^{d_i}$ under the DEM and over all of $\Pset$ under the UAM. Under
\textbf{informed-sincere voting} the same formula is restricted to known projects
$\{q:\sigma_i(q)=1\}$. Under the DEM a voter who knows only a few of their district's projects
still has $\votes$ votes to cast; the leftover votes that cannot stay on known projects are forced
onto unknown ones --- the \emph{noise} we study in Proposition~\ref{prop:noise}.
\end{definition}

\begin{watchout}
Sincere voting is an \emph{assumption about behaviour}, not a prediction of equilibrium. We use it
to isolate the effect of the rule's constraints from the effect of strategic play. When we later
turn to strategy (Propositions~\ref{prop:coord} and~\ref{prop:mob}) we drop it. Be explicit in the
paper about which mode each result assumes --- mixing them up is a common way to state a result more
strongly than it has been proved.
\end{watchout}

\begin{aside}\small
\textbf{A convention we adopt once: divisible votes.} The sincere formula returns real numbers,
which need not be whole votes. Rather than litter the analysis with rounding, we treat votes as
divisible ($x_{i,p}\in\R_{\ge0}$, $\sum_p x_{i,p}=\votes$) throughout the theory; this is standard
and the propositions concern vote \emph{shares}, which rounding does not disturb qualitatively.
Integer ballots ($\votes=6$) re-enter only in the simulation (Section~\ref{sec:sim}), which uses
the real format.
\end{aside}

% ─────────────────────────────────────────────────────────────────────
\subsection{Salience: how visible a district is}
% ─────────────────────────────────────────────────────────────────────

The information function lets us measure something the equity argument will hinge on: how many
voters even know a district exists, project-wise.

\begin{naming}
We call the following quantity the \textbf{salience} of a district. The word is ours (a label for
$\salience_d$ below); the underlying idea --- that visibility, not value, can drive vote totals --- is
the substance.
\end{naming}

\begin{definition}[District salience]
\label{def:salience}
The \textbf{salience} of district $d$ is the fraction of voters informed about at least one of its
projects:
\[
  \salience_d=\frac{\bigl|\{i\in\agents:\exists\,p\in\Pset^d\text{ with }\sigma_i(p)=1\}\bigr|}{n}.
\]
\end{definition}

\begin{watchout}
Salience is built from $\sigma$ (information), \emph{not} from $u$ (value). A glossy city-centre
project can be high-salience and low-value; a vital drainage project in a poor neighbourhood can be
high-value and low-salience. Keeping these two apart is the whole hinge of Proposition~\ref{prop:equity}
and of its reconciliation with Proposition~\ref{prop:trunc}. If you ever find yourself treating
``more votes'' as ``more valued,'' stop --- that is precisely the inference the paper exists to refute.
\end{watchout}

\begin{aside}\small
\textbf{A second convention: compare \emph{scope}, not raw message spaces.} It is tempting to write
$M^{\DEM}\subsetneq M^{\UAM}\subsetneq M^{\MES}$ as if the three rules' message spaces nest. They do
not nest cleanly --- the first two are spaces of \emph{numerical vote allocations}, while the Method
of Equal Shares (Section~\ref{sec:mes}) uses \emph{approval sets}, a different kind of object. What
genuinely nests is the \emph{scope of projects a voter can put weight on}: the DEM confines that
scope to one district; the UAM opens it to the whole city; the Method of Equal Shares additionally
removes the cap on how many projects you may back. We always make the comparison at the level of
scope, never by pretending the message spaces are the same type of set.
\end{aside}

% ═════════════════════════════════════════════════════════════════════
\section{Five ways the DEM distorts outcomes}
\label{sec:distortions}
% ═════════════════════════════════════════════════════════════════════

This is the diagnostic heart of the paper. We prove five propositions, each isolating one way the
district-exclusivity constraint degrades outcomes. The first two hold behaviour fixed at sincere
voting (so they are about the \emph{rule}); the next two are about \emph{strategy} (so we let
voters optimise); the fifth is the twist that stops the paper from being a one-line ``deregulate''
argument.

% ─────────────────────────────────────────────────────────────────────
\subsection{Proposition 1: preference truncation}
% ─────────────────────────────────────────────────────────────────────

The cleanest failure first. A voter who cares about projects in two districts must abandon one of
them at the ballot. We show this costs welfare --- but we have to be careful about \emph{in what
sense}, because the obvious-looking proof is wrong.

\begin{proposition}[Truncation: attainability and welfare]
\label{prop:trunc}
Assume at least two districts and at least one voter with \emph{cross-district preferences}:
there are $p\in\Pset^d$ and $q\in\Pset^{d'}$ with $d\ne d'$ and $u_i(p),u_i(q)>0$. Then:
\begin{enumerate}[label=\textup{(\roman*)},nosep]
  \item \textbf{(Constraint effect.)} For any outcome rule $g$, the DEM's reachable vote profiles
        are contained in the UAM's, $\mathbf{X}^{\DEM}\subseteq\mathbf{X}^{\UAM}$; hence the
        \emph{best attainable} welfare is weakly higher under the UAM.
  \item \textbf{(Truncation.)} Under full information ($\sigma\equiv1$) and sincere voting, the
        DEM vote vector is a strict truncation of the UAM vote vector: it places zero weight on
        every positively-valued project outside the chosen district.
  \item \textbf{(Welfare, conditional.)} Under full information, sincere voting, \emph{and} an
        outcome rule that funds the welfare-best feasible set given the submitted votes,
        $W^{\DEM}\le W^{\UAM}$, with strict inequality whenever a cross-district voter is pivotal
        at the budget cutoff.
\end{enumerate}
\end{proposition}

\begin{proof}
\textbf{(i)} Under the DEM a voter must zero out every project off the chosen district; under the
UAM the only constraint is $\sum_{p\in\Pset}x_{i,p}=\votes$, $x_{i,p}\ge0$. Every DEM profile is a
UAM profile (all weight in one district), so $\mathbf{X}^{\DEM}\subseteq\mathbf{X}^{\UAM}$.
Maximising one and the same objective $W\circ g$ over a weakly larger set cannot lower the maximum.

\smallskip\noindent\textbf{(ii)} With full information, sincere UAM voting gives
$x_{i,p}^{\UAM}=\votes\,u_i(p)/\sum_{q\in\Pset}u_i(q)$ for every $p$, while sincere DEM voting gives
$x_{i,p}^{\DEM}=\votes\,u_i(p)/\sum_{q\in\Pset^{d_i^*}}u_i(q)$ for $p\in\Pset^{d_i^*}$ and $0$
otherwise, where $d_i^*=\argmax_d\sum_{p\in\Pset^d}u_i(p)$ is the district the sincere voter
selects. A cross-district voter has some $q^*\notin\Pset^{d_i^*}$ with $u_i(q^*)>0$; there
$x_{i,q^*}^{\DEM}=0<x_{i,q^*}^{\UAM}$. So the DEM vector is the UAM vector with all out-of-district
weight deleted and the rest renormalised --- a strict truncation.

\smallskip\noindent\textbf{(iii)} A welfare-respecting rule funds the welfare-best feasible set
\emph{consistent with the votes it is handed}. By (ii) the UAM is handed a vector that tracks each
voter's valuation over the whole city; the DEM is handed one that tracks valuation only inside
$d_i^*$. The truncated profile is itself attainable under the UAM (part (i)), and it carries
strictly less valuation information whenever a cross-district voter sits at the funding margin.
Hence the DEM outcome is weakly --- and at a binding pivot strictly --- welfare-dominated.
\end{proof}

\begin{watchout}
Here is the trap, and it is worth dwelling on because the ``natural'' proof is invalid.
One is tempted to argue: ``$\mathbf{X}^{\DEM}\subseteq\mathbf{X}^{\UAM}$, so the UAM optimises over a
bigger set, so $W^{\DEM}\le W^{\UAM}$, full stop.'' \emph{This does not follow.} Sincere voters do
not optimise --- they land on one fixed point of the feasible set --- and the greedy-by-votes funding
rule is \textbf{not monotone} in the vote vector: spreading the same support over more projects can
push a good project \emph{below} the cutoff. So set-containment bounds the \emph{best attainable}
welfare (part (i)); it does \emph{not}, by itself, rank the realised sincere outcomes. That is why
part (iii) carries two hypotheses --- full information \emph{and} a welfare-respecting rule --- and
why we keep the rigorous content (i)--(ii) separate from the welfare claim.

And one phrase to avoid entirely: do \textbf{not} call the DEM outcome ``Pareto-inferior.''
Pareto-inferior means \emph{every} voter is weakly worse off, which is false --- a voter whose
favourite gets funded faster under the DEM's concentration can be strictly better off. The honest
claim is about \emph{utilitarian} welfare, under stated conditions.
\end{watchout}

\begin{intuition}
Strip away the formalism: the DEM forces a voter with three loves to publicly declare only one.
Across a city where people live in one district, work in another, and study in a third, that is a
lot of discarded information. Part (iii) says: when that discarded information would have changed
who sits at the funding margin, total welfare falls. When it would not have mattered, it does not.
\end{intuition}

\begin{aside}\small
This proposition is the ``easy'' one --- the low-hanging diagnosis. It also has a sting in its tail
that we resolve only in Proposition~\ref{prop:equity}: the welfare ranking it states can
\emph{reverse} once information is incomplete. Hold that thought; it is the pivot of the paper.
\end{aside}

\begin{example}[Truncation costs welfare --- with numbers]
\label{ex:trunc}
Two districts, full information, sincere voting; one vote per voter (divisible). District $A$ has a
single project $a$ (cost $10$), district $B$ a single project $b$ (cost $10$), and the budget funds
exactly one. The electorate:
\begin{itemize}[nosep]
  \item $6$ \emph{cross-district} voters with $u(a)=10$, $u(b)=9$;
  \item $5$ \emph{local-$B$} voters with $u(b)=10$, $u(a)=0$.
\end{itemize}
\emph{Under the DEM.} A sincere voter enters the district of higher total utility. For the
cross-district voters that is $A$ ($10>9$), so all six enter $A$ and put their vote on $a$; the five
local voters enter $B$ and back $b$. Totals $V(a)=6$, $V(b)=5$; greedy funds $a$. Welfare
$W^{\DEM}=\sum_i U_i(\{a\})=6(10)+5(0)=60$.
\emph{Under the UAM.} The cross-district voters now split their single vote sincerely, $\tfrac{10}{19}$
to $a$ and $\tfrac{9}{19}$ to $b$. So $V(a)=6\cdot\tfrac{10}{19}\approx3.16$ and
$V(b)=6\cdot\tfrac{9}{19}+5\approx7.84$; greedy funds $b$. Welfare $W^{\UAM}=6(9)+5(10)=104$.
\emph{Read-off.} Project $b$ is in fact the higher-welfare choice ($104$ versus $60$), but the DEM's
district lock-in forced the six cross-district voters to pile all their weight onto $a$ --- funding the
worse project. The truncation cost $104-60=44$ in utilitarian welfare.
\end{example}

\begin{remark}
This is the \emph{well-behaved} case of Proposition~\ref{prop:trunc}: full information, and the greedy
rule happens to fund the welfare-best project under the UAM. Remove full information and the ranking
can flip --- which is precisely Proposition~\ref{prop:equity}. The example is meant to make the mechanism
of part~(iii) tangible, not to suggest the welfare ranking is unconditional.
\end{remark}

% ─────────────────────────────────────────────────────────────────────
\subsection{Proposition 2: noise injection}
% ─────────────────────────────────────────────────────────────────────

Now bring back incomplete information. You enter a district for the one project you know, but you
hold $\votes$ votes and the ballot has many more projects. What happens to the rest of your votes?

\begin{aside}\small
\textbf{The noise model.} Write $I_i^d=\{p\in\Pset^d:\sigma_i(p)=1\}$ for $i$'s informed set in
district $d$ and $\bar I_i^d=\Pset^d\setminus I_i^d$ for the rest. We assume a voter casts informed
votes sincerely and, under social pressure to ``use all six,'' spreads the leftover
$\votes-\votes_i^{\text{inf}}$ votes \emph{uniformly at random} over the unknown projects, where
\[
  \votes_i^{\text{inf}}=\min\!\Bigl(\votes,\ \Bigl\lfloor \votes\cdot
  \tfrac{\sum_{p\in I_i^d}u_i(p)}{\sum_{p\in I_i^d}u_i(p)+\epsilon\,|\bar I_i^d|}\Bigr\rfloor\Bigr),
  \qquad \epsilon>0 .
\]
So $\E[x_{i,p}]=(\votes-\votes_i^{\text{inf}})/|\bar I_i^d|$ for each $p\in\bar I_i^d$. The
parameter $\epsilon$ is the small default weight a voter gives to projects they do not know; the
results hold for any $\epsilon>0$.
\end{aside}

\begin{naming}
We call a vote $x_{i,p}>0$ with $\sigma_i(p)=0$ a \textbf{noise vote}, and the expected total of
such votes on a project its \textbf{noise} $\eta(p)$, as opposed to its \textbf{signal} $S(p)$ from
informed voters. The signal/noise language is our framing for this setting.
\end{naming}

\begin{proposition}[Noise injection]
\label{prop:noise}
Suppose some voter who selected district $d$ is uninformed about a project there. Then for each
$p\in\Pset^d$:
\begin{enumerate}[label=\textup{(\roman*)},nosep]
  \item \textbf{(Decomposition.)} $\E[V(p)]=S(p)+\eta(p)$, with
        $S(p)=\sum_{i:\,d_i=d,\,\sigma_i(p)=1}x_{i,p}^{\text{inf}}$ and $\eta(p)$ the expected
        noise.
  \item \textbf{(Quality-blindness.)}
        $\E[\eta(p)]=\sum_{i:\,d_i=d,\,\sigma_i(p)=0}(\votes-\votes_i^{\text{inf}})/|\bar I_i^d|$,
        in which no voter's $u_j(p)$ appears: the noise on $p$ does not depend on how good $p$ is.
  \item \textbf{(Noise can swamp signal.)} There exist low-signal projects for which
        $\E[\eta(p)]>S(p)$ --- more noise than signal.
\end{enumerate}
\end{proposition}

\begin{proof}
\textbf{(i)} Split the district-$d$ electorate into those informed and those uninformed about $p$;
$V(p)$ is the sum of the two groups' votes, and linearity of expectation gives the decomposition.
\textbf{(ii)} Each uninformed voter contributes
$\E[x_{i,p}]=(\votes-\votes_i^{\text{inf}})/|\bar I_i^d|$, which depends on the mechanism parameter
$\votes$, on $i$'s preferences over \emph{informed} projects (through $\votes_i^{\text{inf}}$), and
on the size of $i$'s unknown set --- but never on $u_j(p)$. Summing preserves this.
\textbf{(iii)} Pick a project with a small informed base (so $S(p)$ is small) about which many
district voters are uninformed (so $\E[\eta(p)]$ collects noise from the whole uninformed pool);
then $\E[\eta(p)]>S(p)$.
\end{proof}

\begin{watchout}
Two claims one is tempted to overstate. First, the quality-blindness in (ii) is true \emph{by
construction} of the uniform-noise model --- present it as a consequence of the modelling
assumption, not as a discovered law of nature. Second, do \textbf{not} claim that the
noise-to-signal ratio is \emph{monotone} in $|\bar I_i^d|/|\Pset^d|$. Lowering how much voters know
raises the leftover noise per voter but spreads it over more unknown projects, so the per-project
effect is ambiguous. The robust, provable claim is the \emph{existence} of noise-swamped projects
(part (iii)); the example below carries the point, and that is all the paper needs.
\end{watchout}

\begin{example}[Durango-scale noise]
\label{ex:noise}
District $A$: $8$ projects, $100$ voters --- $40$ ``FECA'' voters who know only $p_1$, $30$ locals
who know $3$ projects each (including $p_1$), $30$ mobilised voters who each know $1$ project (not
$p_1$).
\emph{Popular $p_1$:} signal $40(4)+30(2)=220$; noise $\approx30(1/7)\approx4.3$ --- about
\textbf{2\%} noise. The obvious winner is read cleanly.
\emph{Small, genuinely valued $p_7$:} signal $5(2)=10$; noise from FECA $40(2/7)\approx11.4$ plus
mobilised $30(1/7)\approx4.3$, total $\approx15.7$ --- about \textbf{61\%} noise. Whether $p_7$ clears
the funding margin is decided more by coin-flips than by anyone's preferences.
\end{example}

\begin{connection}
The policy reading is sharp: the DEM is fine at picking the obvious winners and close to a lottery
on the small, locally-valued projects --- exactly the projects participatory budgeting is supposed to
surface. A rule that let voters simply \emph{not} vote on what they do not know would erase this
noise at the root. That rule is coming in Section~\ref{sec:mes}.
\end{connection}

% ─────────────────────────────────────────────────────────────────────
\subsection{Proposition 3: the coordination game}
% ─────────────────────────────────────────────────────────────────────

Now we let voters be strategic. The district choice is a move in a game, and that game has more
than one equilibrium.

\begin{aside}\small
\textbf{Decomposing the DEM into two stages.} Read the DEM as: first a \emph{strategic} stage
(which district?), then a \emph{sincere} stage (allocate within it). The induced
district-selection game is $\Gamma=\langle\agents,(\D)_i,(\pi_i)_i\rangle$, with payoff
$\pi_i=\E[\,U_i(g^{\DEM})\mid\text{districts},\text{ sincere within-district voting}\,]$. This
split lets us study the coordination incentive in isolation.
\end{aside}

\begin{lemma}[Strategic complementarity in district choice]
\label{lem:complement}
For a voter who values $p^*\in\Pset^d$, the marginal value of joining district $d$ is
\emph{increasing} in the number of others who join when $p^*$ sits near the funding cutoff
(bandwagon), and \emph{decreasing} once $p^*$ is safely funded (diminishing returns).
\end{lemma}

\begin{proof}
Under greedy funding $p^*$ is funded iff $V(p^*)$ ranks high enough city-wide. Near the cutoff,
extra district-$d$ votes compound with the voter's own and raise the chance $p^*$ clears, so
joining is more valuable the more others join. Far above the cutoff, extra votes are wasted on a
project that is already safe, and the voter would do better where their vote is pivotal.
\end{proof}

\begin{proposition}[Multiple equilibria]
\label{prop:coord}
With at least two districts and a block of \emph{mobile} voters --- those with positive utility in
two districts --- large enough to be pivotal, the district-selection game has multiple Nash
equilibria that fund different project sets and yield different welfare.
\end{proposition}

\begin{proof}
Take two districts $A,B$ with committed bases for $p_A,p_B$, neither funded by its base alone, but
either funded if the mobile block joins it.
\emph{Equilibrium $E_A$:} all mobile voters choose $A$; then $p_A$ clears and $p_B$ does not. A lone
defector to $B$ cannot by themselves push $p_B$ over the line and forfeits their contribution to
$p_A$ --- no profitable deviation.
\emph{Equilibrium $E_B$:} symmetric, funding $p_B$ not $p_A$.
The two equilibria fund different sets, and if $\sum_i u_i(p_A)\ne\sum_i u_i(p_B)$ their welfare
differs. Both are Nash, so multiplicity holds.
\end{proof}

\begin{example}[Coordination failure]
\label{ex:coord}
Two districts; a project needs $200$ votes to clear. $A$ has $25$ committed voters
($25\times5=125$ for $p_A$); $B$ has $20$ ($20\times5=100$ for $p_B$); $20$ mobile voters value
$p_A$ at $8$ and $p_B$ at $10$. In $E_A$, $p_A$ gets $125+20(5)=225$ and clears; a switch to $B$
yields $100+5=105<200$. In $E_B$, $p_B$ gets $100+20(5)=200$ and clears. The mobile voters
\emph{prefer} $p_B$ ($10>8$), yet $E_A$ is still an equilibrium: if everyone expects the crowd to
go to $A$, following the crowd is individually rational even though it funds the worse project.
\end{example}

\begin{watchout}
The ``$200$-vote threshold'' in the example is a convenience, not part of the DEM. The real funding
rule is greedy-by-budget. With two projects and budget for one, the rule ``the project with more
votes wins'' reproduces both equilibria ($E_A$: $225>100$; $E_B$: $200>125$), so nothing depends on
positing an exogenous threshold. Say so in the paper rather than smuggling in a threshold.
\end{watchout}

\begin{keyidea}
Multiplicity means the outcome turns on a \emph{focal point}: which district gets talked about,
which group mobilises first, where attention lands. The DEM aggregates \emph{attention} as much as
preference. That is the bridge to the next result.
\end{keyidea}

% ─────────────────────────────────────────────────────────────────────
\subsection{Proposition 4: mobilization advantage}
% ─────────────────────────────────────────────────────────────────────

If coordinating on a district is what wins, then whoever can \emph{make} people coordinate gains
power. A university, an employer, a church --- any body that can say ``everyone vote in district
$A$'' is an equilibrium-selection device.

\begin{naming}
We call such a coordinating body a \textbf{mobilizing institution}, and we call the votes it
diverts from members' own districts \textbf{vote displacement} (written $\Delta_\mathcal{C}$ below).
Both terms are ours; the phenomenon is the familiar one of organised blocs, here given a precise
accounting inside the DEM.
\end{naming}

\noindent Formally, a mobilizing institution is a coalition $\mathcal{C}\subseteq\agents$ with a
\emph{focal project} $p_\mathcal{C}\in\Pset^{d_\mathcal{C}}$, the ability to steer all members to
$d_\mathcal{C}$, and members whose own preferred district $d_i^*$ differs from $d_\mathcal{C}$.

\begin{proposition}[Mobilization and displacement]
\label{prop:mob}
For a mobilizing institution $\mathcal{C}$ with $|\mathcal{C}|=c$:
\begin{enumerate}[label=\textup{(\roman*)},nosep]
  \item it adds exactly $c\votes$ votes to district $d_\mathcal{C}$;
  \item each diverted member inflates $d_\mathcal{C}$ by $\votes$ and deflates their own $d_i^*$ by
        $\votes$;
  \item the net \textbf{vote displacement} is
        $\Delta_\mathcal{C}=\votes\cdot|\{i\in\mathcal{C}:d_i^*\ne d_\mathcal{C}\}|$;
  \item the \textbf{private utility loss to diverted members} is
        \[
          \mathcal{L}(\mathcal{C})=\sum_{i\in\mathcal{C}^*}\bigl[\,U_i(g^*_{d_i^*})-U_i(g^*_{d_\mathcal{C}})\,\bigr]\ \ge\ 0,
        \]
        where $\mathcal{C}^*=\{i\in\mathcal{C}:d_i^*\ne d_\mathcal{C}\}$ and $g^*_d$ is the funded
        set member $i$ would obtain by voting sincerely in district $d$; each term is taken
        conditional on the member's displaced votes being pivotal, and is $0$ otherwise.
\end{enumerate}
\end{proposition}

\begin{proof}
Parts (i)--(iii) are accounting: every voter places exactly $\votes$ votes in one district, so the
DEM is zero-sum across districts and each vote added to $d_\mathcal{C}$ is one withheld from
$d_i^*$. For (iv), absent mobilization member $i$ obtains $U_i(g^*_{d_i^*})$; mobilized, $i$ obtains
$U_i(g^*_{d_\mathcal{C}})$, which already \emph{includes} whatever utility $i$ gets from the focal
project $p_\mathcal{C}$. When the displaced votes are pivotal the move weakly worsens $i$'s own
outcome, so each term is $\ge0$; when they are not pivotal we set the term to $0$. Summing over
$\mathcal{C}^*$ gives the bound.
\end{proof}

\begin{watchout}
Two corrections live here. First, in (iv) compare \emph{full} district utilities:
$U_i(g^*_{d_\mathcal{C}})$ must include the focal project $p_\mathcal{C}$. An earlier version
subtracted only the non-focal projects, dropping $u_i(p_\mathcal{C})$ from the right-hand side and
making the bound sign-ambiguous. Second, name $\mathcal{L}(\mathcal{C})$ honestly: it is a
\textbf{private} loss to the diverted members, \emph{not} a change in system-wide welfare. It says
nothing about non-members, whose own districts now win or lose because of the displacement. Calling
it ``the welfare loss from mobilization'' overclaims.
\end{watchout}

\begin{example}[A university bloc]
A faculty mobilises $c=400$ members to the university's district. Vote displacement is
$\Delta_\mathcal{C}=400\times6=2{,}400$: that many votes are pulled out of members' home districts,
where they can no longer support the projects those members actually value most. The focal project's
total swells by the same amount. None of this is fraud; it is rational play that the DEM rewards.
\end{example}

% ─────────────────────────────────────────────────────────────────────
\subsection{Proposition 5: the equity tradeoff (the pivot)}
% ─────────────────────────────────────────────────────────────────────

So far the DEM looks indefensible. This proposition is what makes the paper interesting: removing
the constraint (the UAM) is \emph{not} a free win, because the constraint, for all its faults,
provides a crude geographic protection that the UAM destroys.

\begin{naming}
We measure that protection with a criterion we call \textbf{$\beta$-geographic equity}. The phrase
is our label for the inequality in Definition~\ref{def:geo}; the idea that no district with genuine
demand should be shut out entirely is, of course, not ours.
\end{naming}

\begin{definition}[$\beta$-geographic equity]
\label{def:geo}
A funded set $S$ satisfies \textbf{$\beta$-geographic equity} if every district with at least one
positively-valued project receives at least a $\beta$-fraction of the per-district equal share:
\[
  \sum_{p\in S\cap\Pset^d}\cost(p)\ \ge\ \beta\cdot\frac{\budget}{m}
  \qquad\text{for every } d \text{ with some } p\in\Pset^d,\ \textstyle\sum_i u_i(p)>0 .
\]
Here $\beta=1$ is equal funding per district and $\beta=0$ is no guarantee at all.
\end{definition}

\begin{proposition}[Equity tradeoff]
\label{prop:equity}
Suppose salience is heterogeneous, $\salience_{d_H}>\salience_{d_L}$ for a high- and a low-salience
district. Then:
\begin{enumerate}[label=\textup{(\roman*)},nosep]
  \item \textbf{(UAM vote ceiling.)} Under sincere UAM voting, every project $p$ obeys
        $V^{\UAM}(p)\le\salience_{d(p)}\,n\votes$; so a low-salience project has a strictly lower
        vote \emph{ceiling}, regardless of its true value.
  \item \textbf{(DEM shelter.)} Under the DEM, the exclusivity constraint segments competition:
        committed residents of $d_L$ supply a pool of votes to $d_L$'s projects that cannot leak
        away to $d_H$.
  \item \textbf{(The gap.)} There are parameter values where the DEM achieves $\beta>0$ while the
        UAM yields $\beta=0$ --- the low-salience district is shut out completely.
\end{enumerate}
\end{proposition}

\begin{proof}
\textbf{(i)} Under sincere UAM voting a voter puts zero votes on projects they do not know, so
$V^{\UAM}(p)=\sum_{i:\sigma_i(p)=1}x_{i,p}\le|\{i:\sigma_i(p)=1\}|\,\votes\le\salience_{d(p)}\,n\votes$,
the last step being the definition of salience.
\textbf{(ii)} If $n_{d_L}$ residents select $d_L$ they direct $\votes\,n_{d_L}$ votes onto $d_L$'s
projects, sealed off from competition with the high-salience centre.
\textbf{(iii)} Take $\salience_{d_H}=0.8$, $\salience_{d_L}=0.1$, $n=1000$, $\votes=6$. Under the
UAM, $d_H$'s projects can reach $4{,}800$ votes and $d_L$'s at most $600$; if the top funded
projects are all high-salience, $d_L$ gets nothing, $\beta=0$. Under the DEM, $100$ committed $d_L$
residents supply $600$ concentrated votes over (say) three local projects --- enough to fund the top
one, so $\beta>0$.
\end{proof}

\begin{watchout}
\textbf{This is where Proposition~1 and Proposition~5 must be reconciled, or a referee will catch a
contradiction.} Proposition~1 says (under its hypotheses) $W^{\DEM}\le W^{\UAM}$; Proposition~5
builds a case where the UAM funds nothing in a district whose projects are genuinely valued ---
a welfare \emph{gain} for the DEM. Both cannot hold unconditionally. The resolution is the
full-information hypothesis in Proposition~\ref{prop:trunc}(iii): the reversal here is driven by
\emph{salience} (an information object), not by \emph{utility}. Under the UAM a widely-\emph{known}
low-value project can outpoll a narrowly-known high-value one --- a channel that exists only when
$\sigma\not\equiv1$, and which full information switches off. So the two propositions are
complementary: Proposition~1 is the \emph{constraint} effect (full information), Proposition~5 the
\emph{information} effect (salience asymmetry). Together they say: deregulation alone is not the
answer, which is the entire motivation for the next chapter.
\end{watchout}

\begin{aside}\small
\textbf{One open modelling choice to settle on the page.} Definition~\ref{def:geo} benchmarks
against the \emph{equal-per-district} share $\budget/m$. The guarantee we prove for the Method of
Equal Shares (Proposition~\ref{prop:mesq}) is instead \emph{population-proportional},
$\beta_d\budget$ with $\beta_d=n_d/n$. These are different fairness ideals, and they do not coincide
--- a small marginalised district fares better under equal-per-district than under
population-proportional. Pick one benchmark for the paper and use it consistently; the cleanest
choice, absent other weighting, is population-proportional throughout.
\end{aside}

% ═════════════════════════════════════════════════════════════════════
\section{A better rule: the Method of Equal Shares}
\label{sec:mes}
% ═════════════════════════════════════════════════════════════════════

Propositions 1--4 indict the DEM; Proposition~5 warns that mere deregulation (the UAM) trades one
failure for another. We need a rule that expresses cross-district preferences \emph{and} protects
small districts by construction. Such a rule already exists in the social-choice literature.

\begin{aside}\small
\textbf{This one is not our coinage.} The \textbf{Method of Equal Shares} (MES) and the fairness
axiom it satisfies, \textbf{Extended Justified Representation} (EJR), are established results of
\citet{peters2021}; likewise ``cohesive group'' is standard terminology there. We adopt the names
as-is and cite the source. (Contrast the violet boxes earlier, which flagged labels we invented.)
\end{aside}

% ─────────────────────────────────────────────────────────────────────
\subsection{The idea, then the algorithm}
% ─────────────────────────────────────────────────────────────────────

\begin{keyidea}
Treat the budget as if it were \emph{individually owned}. Give every voter a virtual share
$\budget/n$. A project gets funded when the voters who back it can jointly pay its cost out of
their remaining shares, charging supporters as \emph{equally} as the budgets allow. The
consequence is proportional: any group of $k$ voters controls $k\cdot\budget/n$ that no majority
can spend on its behalf. A marginalised district's residents thus carry their own protected
purse --- no district lock-in required.
\end{keyidea}

\begin{definition}[Method of Equal Shares]
\label{def:mes}
Voters submit \textbf{approval ballots}: $A_i\subseteq\Pset$ is the set of projects $i$ supports
(no district constraint). Initialise each voter's virtual budget at $b_i=\budget/n$. Then iterate:
\begin{enumerate}[label=\textbf{\arabic*.},nosep]
  \item For each unfunded project $p$, let $N(p)=\{i:p\in A_i,\ b_i>0\}$ be its active supporters.
  \item For each such $p$, compute the \textbf{per-supporter cost cap}
        \[
          \mesrho(p)=\min\Bigl\{\rho\ge0:\ \textstyle\sum_{i\in N(p)}\min(b_i,\rho)\ \ge\ \cost(p)\Bigr\}.
        \]
  \item If no project has a finite $\mesrho$, \textbf{stop}.
  \item Otherwise fund $p^*=\argmin_p \mesrho(p)$; each supporter $i$ pays $\min(b_i,\mesrho(p^*))$,
        and $b_i$ drops by that amount.
  \item Remove $p^*$ and repeat.
\end{enumerate}
\textbf{Completion.} If budget remains when MES halts, a completion rule (typically greedy by
number of approvals) spends the rest.
\end{definition}

\begin{intuition}
$\mesrho(p)$ is the smallest \emph{equal} charge per supporter that covers $\cost(p)$, with
supporters too poor to pay $\rho$ capped at whatever they have left and the rest made up by the
richer backers. A small $\mesrho(p)$ means ``cheap per supporter'' --- either the project is cheap
or it has many well-funded backers. Funding the smallest-$\mesrho$ project first rewards broad,
affordable support and quietly screens out expensive projects held by a narrow group.
\end{intuition}

\begin{remark}[Why $\mesrho(p)$ is well-defined]
Let $h(\rho)=\sum_{i\in N(p)}\min(b_i,\rho)$. It is continuous and non-decreasing in $\rho$, with
$h(0)=0$ and $h(\rho)\uparrow\sum_{i\in N(p)}b_i$. So a finite $\mesrho(p)$ exists \emph{iff}
$\sum_{i\in N(p)}b_i\ge\cost(p)$; otherwise the supporters simply cannot afford $p$ and it is
skipped --- exactly the stopping condition in step~3.
\end{remark}

% ─────────────────────────────────────────────────────────────────────
\subsection{A walkthrough}
% ─────────────────────────────────────────────────────────────────────

\begin{example}[The Method of Equal Shares, step by step]
\label{ex:mes}
$\budget=100$, $n=4$, so $b_i=25$ each. Approvals: $p_1$ (cost $40$) by voters $1,2,3$; $p_2$
(cost $30$) by $3,4$; $p_3$ (cost $50$) by $1,2$; $p_4$ (cost $20$) by $4$.

\smallskip\noindent\textbf{Round 1.}
$\mesrho(p_1)=40/3\approx13.3$, $\mesrho(p_2)=15$, $\mesrho(p_3)=25$, $\mesrho(p_4)=20$.
Smallest is $p_1$: \textbf{fund $p_1$}; voters $1,2,3$ each pay $13.3$, leaving
$b_1=b_2=b_3=11.67$, $b_4=25$.

\smallskip\noindent\textbf{Round 2.}
$p_2$: supporters $3,4$ hold $11.67$ and $25$; solving $11.67+\rho=30$ gives $\mesrho(p_2)=18.33$.
$p_4$: supporter $4$ holds $25$, so $\mesrho(p_4)=20$.
$p_3$: supporters $1,2$ hold $11.67$ each, total $23.34<50$ --- not affordable.
The smallest is $p_2$ (since $18.33<20$): \textbf{fund $p_2$}; voter $3$ pays $11.67$ (to $b_3=0$),
voter $4$ pays $18.33$ (to $b_4=6.67$).

\smallskip\noindent\textbf{Round 3.}
$p_4$: only voter $4$ remains, with $6.67<20$ --- not affordable. $p_3$: still $23.34<50$. MES
\textbf{stops} with $\{p_1,p_2\}$, cost $70$, leaving $30$.

\smallskip\noindent\textbf{Completion.} The leftover $30$ funds $p_4$ (cost $20\le30$); $p_3$
(cost $50$) cannot be afforded. \textbf{Final outcome: $\{p_1,p_2,p_4\}$}, cost $90$.

\smallskip\noindent\emph{What to notice.} Both marginal-district projects are funded with no
geographic rule at all --- $p_2$ \emph{inside} MES (its two backers share the cost at a lower cap
than $p_4$'s lone backer) and $p_4$ via completion. The expensive, narrowly-held $p_3$ is screened
out once its backers have spent on $p_1$.
\end{example}

\begin{watchout}
A subtle step that is easy to botch: in Round 2 the project to fund is $p_2$, not $p_4$, because
$\mesrho(p_2)=18.33<20=\mesrho(p_4)$. (An earlier version of this example funded $p_4$ here; the
final set happens to be the same, but the \emph{path} --- and therefore the explanation --- is
different: $p_2$ is funded \emph{inside} MES, $p_4$ only in completion.) Always recompute every
$\mesrho$ from the \emph{current} budgets each round; intuitions about ``the cheapest project''
formed in Round~1 do not survive the budget updates.
\end{watchout}

% ─────────────────────────────────────────────────────────────────────
\subsection{Why it protects minorities: EJR and geographic equity}
% ─────────────────────────────────────────────────────────────────────

The protection is not a happy accident of the example; it is a theorem.

\begin{definition}[Cohesive group; Extended Justified Representation]
\label{def:ejr}
A set of voters $S$ is \textbf{$T$-cohesive} for a project set $T$ if every member approves every
project in $T$ ($T\subseteq A_i$ for all $i\in S$). An outcome $W$ satisfies \textbf{EJR} if for
every $T$-cohesive group $S$ whose shares can afford $T$ --- that is,
$|S|\,\budget/n\ge\sum_{p\in T}\cost(p)$ --- at least one member $i\in S$ is represented to that
level: $\sum_{p\in W\cap A_i}\cost(p)\ge\sum_{p\in T}\cost(p)$.
\end{definition}

\begin{intuition}
EJR is a floor on how thoroughly any large-enough, agreed-enough group can be ignored. If a
group's own shares could have bought a bundle $T$, the rule must leave at least one member as well
represented as $T$. It is the precise sense in which a minority ``cannot be outvoted into nothing.''
\end{intuition}

\begin{theorem}[The Method of Equal Shares satisfies EJR; \citealp{peters2021}]
\label{thm:ejr}
The outcome of the Method of Equal Shares satisfies EJR.
\end{theorem}

\begin{proof}[Proof idea (the full proof is \citealp{peters2021}; we cite it rather than reprove)]
Suppose a $T$-cohesive group $S$ that can afford $T$ ends up under-represented. Track $S$'s
collective budget: it begins at $|S|\,\budget/n$ and only falls when MES funds projects that
$S$-members approve. If some project of $T$ is ever funded, $S$-members paid into it and it counts
toward their representation. If no project of $T$ is ever funded, then at every round a cheaper
(smaller-$\mesrho$) project that $S$-members also approve was chosen instead --- spending their
shares on other approved projects, which also count. Either way the collective share cannot drain
to zero without some member reaching representation $\ge\sum_{p\in T}\cost(p)$, contradicting
under-representation. The complete argument is an induction over the sequence of selections.
\end{proof}

\begin{example}[EJR in action: a minority that cannot be outvoted away]
\label{ex:ejr}
$n=10$ voters, $\budget=10$, so $b_i=1$. A majority of $6$ voters approve three cheap, popular
projects $m_1,m_2,m_3$ (each cost $3$); a minority of $4$ voters approve only one project $q$
(cost $4$). The minority's shares total $4\cdot1=4=\cost(q)$, so they are a $\{q\}$-cohesive group
that can just afford $q$ --- EJR \emph{must} represent them.

\emph{What the Method of Equal Shares does.} Round 1:
$\mesrho(m_1)=\mesrho(m_2)=\mesrho(m_3)=3/6=0.5$ and $\mesrho(q)=4/4=1$; fund $m_1$, the six majority
voters drop to $0.5$. Round 2: each remaining $m_j$ still has $\mesrho=0.5$ (six voters at $0.5$ cover
$3$ exactly); fund $m_2$, the majority drop to $0$. Round 3: $m_3$'s backers are now broke, so it is
unaffordable; $q$'s four backers still hold $1$ each, $\mesrho(q)=1$; fund $q$. Outcome
$\{m_1,m_2,q\}$, cost $10$. The minority got $q$.

\emph{Contrast with a vote-count rule.} Greedy by number of approvals would fund $m_1,m_2,m_3$ (each
with $6$ approvals, total cost $9$), leaving only $1$ --- not enough for $q$ ($4$ approvals, cost $4$).
The minority is shut out. EJR is exactly the property that forbids this: because the minority could
pay for $q$ out of their own shares, the Method of Equal Shares guarantees they are not outvoted into
nothing. This is the engine behind the geographic guarantee we prove next.
\end{example}

\noindent Now translate proportionality into geography. A district's residents who only care about
local projects form exactly the kind of cohesive group EJR protects.

\begin{proposition}[The Method of Equal Shares implies geographic equity]
\label{prop:mesq}
Let district $d$ have $n_d$ \textbf{committed residents} (voters with $A_i\subseteq\Pset^d$).
Suppose there is a \textbf{commonly-approved} local bundle
$T_d\subseteq\bigcap_{i\,\text{resident}}A_i$ --- one that \emph{every} resident approves --- with
$\sum_{p\in T_d}\cost(p)\le n_d\,\budget/n$. Then some resident is represented to the level of
$T_d$, and hence
\[
  \sum_{p\in W\cap\Pset^d}\cost(p)\ \ge\ \sum_{p\in T_d}\cost(p).
\]
Writing $C_d^\star$ for the most expensive such commonly-approved affordable bundle, district $d$ is
guaranteed funding $C_d^\star$; and if the residents commonly approve local projects worth at least
$\beta_d\budget$ (with $\beta_d=n_d/n$), the guarantee is the full population share $\beta_d\budget$.
\end{proposition}

\begin{proof}
By hypothesis $T_d\subseteq A_i$ for all $n_d$ residents, so they form a $T_d$-cohesive group of
size $n_d$ with $n_d\,\budget/n\ge\sum_{p\in T_d}\cost(p)$. EJR (Theorem~\ref{thm:ejr}) yields a
resident $i$ with $\sum_{p\in W\cap A_i}\cost(p)\ge\sum_{p\in T_d}\cost(p)$. Since
$A_i\subseteq\Pset^d$, we have $W\cap A_i\subseteq W\cap\Pset^d$, so
$\sum_{p\in W\cap\Pset^d}\cost(p)\ge\sum_{p\in T_d}\cost(p)$. Taking the maximal commonly-approved
affordable bundle gives $C_d^\star$.
\end{proof}

\begin{watchout}
\textbf{Two hypotheses are doing the work, not one.} ``Committed'' ($A_i\subseteq\Pset^d$) only
guarantees that whatever representation EJR delivers \emph{lands inside} $d$. EJR additionally needs
a \textbf{common core} $T_d$ that \emph{all} $n_d$ residents approve --- being committed does not make
residents agree on \emph{which} local projects to fund. If their local tastes are heterogeneous, the
full $\beta_d\budget$ floor can shrink toward the best cohesive sub-coalition's $|S|\,\budget/n$;
it stays strictly positive, hence still better than the UAM's $\beta=0$, but it is not automatically
the population share. State the floor against the common-core bundle, or argue (empirically, in the
simulation) that Durango's districts have a sizeable project most residents back.
\end{watchout}

\begin{connection}
The bottom line for Durango: a district that is, say, 5\% of the electorate and rallies behind a
local project is guaranteed roughly 5\% of the budget for it --- whatever the city centre's
visibility, whatever a university campaign does, whatever voters know about other districts. The
protection is structural, not behavioural, which is exactly what Proposition~5 said the DEM's
protection was \emph{not}.
\end{connection}

% ─────────────────────────────────────────────────────────────────────
\subsection{The Method of Equal Shares undoes the four distortions}
% ─────────────────────────────────────────────────────────────────────

Each of the next four propositions pairs with one earlier failure. Three are nearly definitional;
state them briefly and move on.

\begin{proposition}[No truncation]
\label{prop:mes-trunc}
Under the Method of Equal Shares a ballot $A_i\subseteq\Pset$ may contain projects from any
districts; the reachable \emph{scope} of supported projects strictly contains the DEM's (it removes
both the single-district restriction and the cap on how many projects one may back).
\end{proposition}
\begin{proof}
Immediate from Definition~\ref{def:mes}: there is no district-selection step, so any subset of
$\Pset$ --- in particular any cross-district one --- is a legal ballot. (We phrase the comparison as
scope, per the convention in Section~\ref{sec:model}, not as a nesting of message spaces.)
\end{proof}

\begin{proposition}[No forced noise]
\label{prop:mes-noise}
A voter who does not know project $p$ ($\sigma_i(p)=0$) simply leaves it out of $A_i$. The rule
never demands a minimum ballot size, so no noise votes arise; an under-informed voter's unspent
share flows to completion as unused entitlement, not as misdirected support.
\end{proposition}
\begin{proof}
The DEM forced all $\votes$ votes inside one district, unknown projects included --- the source of
noise. The Method of Equal Shares imposes no such floor, so $A_i$ can be exactly the known, valued
set and every recorded approval is informative.
\end{proof}

\begin{proposition}[No district-selection game]
\label{prop:mes-coord}
There is no district-selection stage, so the coordination game of Proposition~\ref{prop:coord} does
not arise; the residual strategic interaction is project-by-project rather than a coarse,
focal-point-driven district choice.
\end{proposition}
\begin{proof}
The multiplicity in Proposition~\ref{prop:coord} came from the bandwagon complementarity of the
\emph{district} choice. Here a voter's payoff from approving $p$ depends on $p$'s other supporters
(through $\mesrho(p)$) and on their own remaining budget --- a finer interaction with no ``which
district?'' focal point to herd on.
\end{proof}

\begin{watchout}
Proposition~\ref{prop:mes-coord} is a \emph{plausibility} claim --- ``no district stage, so no
bandwagon multiplicity'' --- not an equilibrium-uniqueness theorem. The Method of Equal Shares is
still manipulable (Proposition~\ref{prop:sp} below). Claim only that the DEM's specific coordination
pathology is gone, not that all strategic behaviour vanishes.
\end{watchout}

\begin{proposition}[Bounded mobilization]
\label{prop:mes-mob}
A coalition $\mathcal{C}$ with $|\mathcal{C}|=c$ can fund a focal project $p_\mathcal{C}$
\emph{by itself} only if $\cost(p_\mathcal{C})\le c\,\budget/n$. It cannot touch non-members'
shares, so the DEM's vote displacement is gone ($\Delta=0$).
\end{proposition}
\begin{proof}
Funding $p_\mathcal{C}$ requires its backers to cover the cost from their own budgets; the
coalition's maximum contribution is $\sum_{i\in\mathcal{C}}b_i\le c\,\budget/n$. Non-members keep
their full shares and their EJR-protected representation, and there are no districts to drain.
\end{proof}

\begin{example}[The university bloc, revisited]
A faculty mobilises $200$ of $n=2{,}000$ voters; $\budget=1{,}000{,}000$ pesos; the focal project
costs $120{,}000$. Under the Method of Equal Shares the bloc commands only $200\times500=100{,}000$
--- short of $120{,}000$, so it needs at least $40$ genuine outside supporters. Under the DEM the same
$200$ concentrate $1{,}200$ votes, plausibly capturing the district while stripping $1{,}200$ votes
from home districts. The contrast is the whole point: the new rule makes the bloc \emph{earn} broad
support rather than manufacture it by relocation.
\end{example}

% ─────────────────────────────────────────────────────────────────────
\subsection{An honest limitation: it is not strategy-proof}
% ─────────────────────────────────────────────────────────────────────

A credible paper states what its proposal does \emph{not} fix. The Method of Equal Shares is not
strategy-proof, and we show this with a concrete, robust example.

\begin{proposition}[The Method of Equal Shares is not DSIC]
\label{prop:sp}
There are preference profiles where a voter strictly gains by submitting a ballot other than their
sincere one.
\end{proposition}

\begin{proof}
Take $n=4$, $b_i=1$ each, two projects $p_1,p_2$ each costing $2$. Sincere ballots:
$V_1=\{p_1,p_2\}$, $V_2=V_3=\{p_1\}$, $V_4=\{p_2\}$. Then $\mesrho(p_1)=2/3<1=\mesrho(p_2)$, so MES
funds $p_1$, charging voters $1,2,3$ each $2/3$. Afterward $p_2$'s backers $1,4$ hold
$1/3+1=4/3<2$, so $p_2$ fails: the outcome is $\{p_1\}$.

Now voter $1$ deviates to $A_1'=\{p_2\}$ --- hiding her support for $p_1$. Now $\mesrho(p_1)=1$
(backers $2,3$) and $\mesrho(p_2)=1$ (backers $1,4$). Whichever is funded first at cap $1$, the
other becomes affordable immediately after, so the outcome is $\{p_1,p_2\}$ --- \emph{under either
tie-break}. Voter $1$ moves from $u_1(p_1)$ to $u_1(p_1)+u_1(p_2)$: a strict gain, achieved by
free-riding on voters $2,3$'s support for $p_1$ and saving her own budget for $p_2$.
\end{proof}

\begin{watchout}
Two things to get right when you write this up. First, the example is \textbf{tie-break robust}:
both resolutions of the $\mesrho=1$ tie give $\{p_1,p_2\}$, so do not hedge with ``say $p_2$ wins''
and do not bother perturbing the costs. Second, attribute the impossibility correctly. The reason no
fix is available is \emph{not} Gibbard--Satterthwaite (whose unrestricted-domain hypothesis fails
here, as we stressed in Section~\ref{sec:what}); it is the domain-appropriate result of
\citet{peters2021} that \emph{no rule satisfying EJR can be strategy-proof}. Present
Proposition~\ref{prop:sp} as a concrete instance of that result, and frame manipulability as the
unavoidable price of the very proportionality that buys us geographic equity.
\end{watchout}

\begin{aside}\small
\textbf{Why this is not catastrophic in practice.} Three things keep the manipulability of
Proposition~\ref{prop:sp} from undermining the proposal. \textbf{(1)} The deviation requires detailed
knowledge of how everyone else will vote --- which project will be ``cheap'' enough to free-ride on ---
and in a large electorate that information simply is not available. \textbf{(2)} It is risky: if
enough voters hide their support for the popular project, it may fail to be funded at all, and the
manipulator loses it outright. \textbf{(3)} It is provably unavoidable for any proportional rule
(\citealp{peters2021}), so the Method of Equal Shares is already as good as one can do on this axis.
None of this makes the rule strategy-proof; together they make the residual manipulability low-stakes.
\end{aside}

% ─────────────────────────────────────────────────────────────────────
\subsection{Two more honest limitations: the completion gap and ballot complexity}
% ─────────────────────────────────────────────────────────────────────

Strategy-proofness is not the only thing the Method of Equal Shares gives up. Two practical
limitations deserve their own discussion, because a referee --- or a skeptical municipal official ---
will raise them, and the paper is stronger for meeting them head-on.

\medskip\noindent\textbf{The completion gap.} The rule routinely fails to spend the whole budget in
its proportional phase. Once every remaining project is unaffordable to its backers, MES stops, and a
\emph{completion} step (greedy by approvals) places the rest --- carrying \emph{none} of the EJR
guarantee. The size of that gap is not a footnote: it depends on how project costs compare to the
per-voter share $\budget/n$. Many small projects let the proportional phase reach far; a few expensive
projects leave most of the budget to completion. The next example makes the worry concrete.

\begin{example}[When completion does most of the work]
\label{ex:completion}
$n=10$, $\budget=100$, so $b_i=10$. Projects: $p_1$ (cost $30$, backed by voters $1$--$6$); $p_2$
(cost $60$, backed by voters $5$--$10$); $p_3$ (cost $50$, backed by voters $1$--$5$).
\emph{Round 1:} $\mesrho(p_1)=30/6=5$, $\mesrho(p_2)=60/6=10$, $\mesrho(p_3)=50/5=10$; fund $p_1$, its
six backers drop to $5$.
\emph{Round 2:} $p_2$'s backers can muster at most $2(5)+4(10)=50<60$; $p_3$'s can muster $5(5)=25<50$.
\textbf{MES stops}, having spent just $30$ of $100$. Completion then funds $p_2$ (cost $60$, more
approvals) out of the leftover, ending at $\{p_1,p_2\}$.
\emph{The point:} two-thirds of the spent budget ($60$ of $90$) was placed by \emph{completion}, which
has no proportionality guarantee. With projects expensive relative to $\budget/n$, the part of the rule
that earns our equity theorem barely gets to act.
\end{example}

\begin{connection}
The design implication for Durango is concrete: keep project cost tiers small relative to the
per-voter share $\budget/n$, so the proportional phase --- the part carrying the equity guarantee ---
does as much of the allocation as possible. \citet{peters2021} discuss several completion methods; the
paper should name the one it uses and report what fraction of the budget that completion step ends up
placing.
\end{connection}

\medskip\noindent\textbf{Ballot complexity.} Under the DEM a voter's task is small: pick a district,
split six votes. Under the Method of Equal Shares the voter is asked to consider \emph{every} project,
city-wide, and decide which to approve. With fifty or more projects that is a genuine cognitive load,
and it interacts with the information problem of Section~\ref{sec:distortions}: a voter who knows few
projects submits a thin ballot, and the budget share they leave unspent leaks to completion. The
remedy is interface design, not a change to the rule --- present projects grouped by district with short
descriptions and costs, let voters approve as many or as few as they wish, and show a running ``budget
impact'' indicator so a voter can see how much of their share they have committed. That is
informationally \emph{richer} than the DEM, which gives no sense of opportunity cost at all, while
staying tractable.

% ═════════════════════════════════════════════════════════════════════
\section{The three rules side by side}
\label{sec:comparison}
% ═════════════════════════════════════════════════════════════════════

We can now assemble the diagnosis and the remedy into one comparison, then watch all three rules
run on a single realistic profile.

\begin{theorem}[Mechanism comparison]
\label{thm:comp}
For a participatory-budgeting environment with heterogeneous salience and voters who have
cross-district preferences and incomplete information, the three rules compare as follows:
\begin{center}\small\renewcommand{\arraystretch}{1.2}
\begin{tabular}{p{6.6cm}ccc}
\toprule
\textbf{Property} & \textbf{DEM} & \textbf{UAM} & \textbf{MES}\\
\midrule
Full preference expression (no truncation) & $\times$ & $\checkmark$ & $\checkmark$\\
No forced noise injection & $\times$ & $\checkmark$ & $\checkmark$\\
No coordination game & $\times$ & $\checkmark$ & $\checkmark$\\
Resistant to mobilization capture & $\times$ & Partial & $\checkmark$\\
$\beta$-geographic equity & Partial & $\times$ & $\checkmark$ (via EJR)\\
Strategy-proof (DSIC) & $\times$ & $\times$ & $\times$\\
Proportional representation & $\times$ & $\times$ & $\checkmark$\\
Simple ballot & $\checkmark$ & $\checkmark$ & Partial\\
\bottomrule
\end{tabular}
\end{center}
Moreover: \textbf{(i)} no rule satisfying EJR is strategy-proof (\citealp{peters2021}), so the
Method of Equal Shares' manipulability is inherent, not a design flaw; \textbf{(ii)} its geographic
equity is \emph{constructive} (it follows from EJR, proportionally to population), whereas the DEM's
is \emph{accidental} (it rests on voters self-selecting into home districts, and is undone by
mobilization); \textbf{(iii)} under full information and sincere voting, the Method of Equal Shares
weakly dominates the DEM on \emph{expected utilitarian welfare}.
\end{theorem}

\begin{proof}[Proof (assembly)]
Each table cell is a proposition already proved. \emph{DEM column:} Propositions~\ref{prop:trunc}
(truncation), \ref{prop:noise} (noise), \ref{prop:coord} (coordination), \ref{prop:mob}
(mobilization), \ref{prop:equity} (only partial, accidental equity); not DSIC because the district
stage is manipulable. \emph{UAM column:} no truncation/noise/coordination by construction; only
partial mobilization resistance (a bloc may concentrate votes but cannot stop others voting);
$\times$ equity by Proposition~\ref{prop:equity}; not DSIC (votes can be strategically concentrated).
\emph{MES column:} Propositions~\ref{prop:mes-trunc}--\ref{prop:mes-mob} and
Theorem~\ref{thm:ejr}/Proposition~\ref{prop:mesq}; not DSIC by Proposition~\ref{prop:sp}. Claim (i)
is the impossibility of \citet{peters2021}; (ii) contrasts a structural budget-accounting guarantee
with a behavioural one; (iii) holds because the Method of Equal Shares can reproduce any DEM support
pattern and, under full information, its selection improves expected utilitarian welfare.
\end{proof}

\begin{watchout}
Claim (iii) is \emph{not} Pareto-dominance. (An earlier statement said ``weakly Pareto-dominates''
--- too strong.) A voter whose favourite is a high-salience central project that the DEM's
concentration funds quickly can be \emph{worse} off under the Method of Equal Shares. State only
\emph{expected utilitarian} welfare dominance, and note that even that inherits the
full-information caveat of Proposition~\ref{prop:trunc}. The genuinely unconditional advantages are
the structural rows of the table --- lead with those, not with a welfare number.
\end{watchout}

% ─────────────────────────────────────────────────────────────────────
\subsection{A full worked comparison at Durango scale}
% ─────────────────────────────────────────────────────────────────────

\begin{example}[$n=300$, three districts, all three rules]
\label{ex:300}
$\budget=300$, $\votes=6$ (for DEM/UAM), $b_i=1$ (for MES). Districts and projects:
\begin{center}\small
\begin{tabular}{lccc}
\toprule
 & District $A$ (centre) & District $B$ (residential) & District $C$ (marginalised)\\
\midrule
Committed residents & 100 & 120 & 80\\
Salience $\salience_d$ & 0.90 & 0.50 & 0.15\\
Projects & $p_1$ (120), $p_2$ (80) & $p_3$ (100) & $p_4$ (60), $p_5$ (40)\\
\bottomrule
\end{tabular}
\end{center}
Voter blocks: $100$ in $A$ approve $\{p_1,p_2\}$; $80$ mobile in $B$ approve $\{p_1,p_3\}$; $40$
local in $B$ approve $\{p_3\}$; $60$ in $C$ approve $\{p_4,p_5\}$; $20$ ``connected'' in $C$ approve
$\{p_1,p_4\}$.

\medskip\noindent\textbf{Under the DEM} (with bandwagoning: $50$ of the mobile $B$-voters and the
$20$ connected $C$-voters defect to $A$ for $p_1$):
$V(p_1)=100(4)+50(6)+20(6)=820$, $V(p_3)=30(6)+40(6)=420$, $V(p_4)=60(4)=240$, $V(p_2)=200$,
$V(p_5)=120$. Greedy funds $\{p_1,p_3,p_4\}$ (cost $280$). District $C$ gets $p_4$ --- but $70$ voters
were pulled out of their home districts to get there.

\medskip\noindent\textbf{Under the UAM} (sincere across known projects): roughly $p_1\approx600$,
$p_3\approx360$, $p_2\approx200$, $p_4\approx160$, $p_5\approx60$. Greedy funds $\{p_1,p_2,p_3\}$
(cost $300$). \textbf{District $C$ gets nothing}: $\beta=0$.

\medskip\noindent\textbf{Under the Method of Equal Shares} ($b_i=1$):
\emph{Round 1:} $\mesrho(p_1)=120/200=0.60$ is smallest; fund $p_1$, dropping its $200$ backers to
$0.40$.
\emph{Round 2:} $p_2$ ($40<80$) and $p_3$ ($80(0.40)+40(1)=72<100$) are unaffordable; $p_4$ needs
$\mesrho\approx0.87$; $p_5$ has $\mesrho=40/60\approx0.667$, the smallest --- fund $p_5$, dropping the
$60$ $C$-residents to $0.333$.
\emph{Round 3:} nothing is affordable ($p_4$: $60(0.333)+20(0.40)=28<60$). MES stops with
$\{p_1,p_5\}$ (cost $160$); completion funds $p_3$ (cost $100$), and the remaining $40$ cannot afford
$p_4$. \textbf{Outcome $\{p_1,p_3,p_5\}$}: district $C$ is funded out of \emph{its own residents'
shares}, with no district lock-in, no noise, and no displacement.

\medskip
\begin{center}\small\renewcommand{\arraystretch}{1.2}
\begin{tabular}{lccc}
\toprule
 & DEM & UAM & MES\\
\midrule
District $C$ funded & $p_4$ (partial) & \textbf{nothing} & $p_5$\\
Noise votes & yes & no & no\\
Coordination needed & yes & no & no\\
Vote displacement & yes ($70$) & no & no\\
\bottomrule
\end{tabular}
\end{center}
\end{example}

% ═════════════════════════════════════════════════════════════════════
\section{From theory to simulation}
\label{sec:sim}
% ═════════════════════════════════════════════════════════════════════

The empirical part of this paper is a \textbf{simulation study}, not an analysis of municipal
records. We do not yet have ballot-level data from the \emph{Municipio de Durango}; obtaining it is
deliberately left to follow-up work. The simulation is what we can build now, and it doubles as the
artefact we put in front of the municipality --- a concrete demonstration that might persuade them to
share data for a later study.

\begin{aside}\small
\textbf{Scope, stated plainly.} This paper $=$ theory (Sections~\ref{sec:what}--\ref{sec:comparison})
$+$ simulation (this section). Real-data estimation, and any need-based reweighting of shares, are
explicitly future work. Keeping that boundary clear protects the paper from promising more than it
delivers.
\end{aside}

% ─────────────────────────────────────────────────────────────────────
\subsection{What the simulation establishes}
% ─────────────────────────────────────────────────────────────────────

On a preference profile --- synthetic now, real later if data arrives --- run all three rules on
identical inputs and compare the funded sets on two precise axes: welfare and geographic equity.

\begin{definition}[Welfare measures]
\label{def:welfare}
When the profile is synthetic the utilities $u_i$ are known, so report \textbf{true utilitarian
welfare} $W^M=\sum_i U_i(g^M)=\sum_i\sum_{p\in g^M}u_i(p)$ for each rule $M$. When only ballots are
observed (no $u_i$), fall back to the \textbf{vote-total proxy}
$\widehat W^{M}=\sum_{p\in g^M}V(p)$, reported alongside the \textbf{signal-only proxy}
$\widehat W^{M}_{\text{sig}}=\sum_{p\in g^M}S(p)$ (informed votes only).
\end{definition}

\begin{watchout}
Do not conflate $W$ and $\widehat W$. The vote-total proxy \emph{builds in} the very noise and
salience distortions the paper is criticising --- a project funded mostly by noise scores as ``high
welfare'' under $\widehat W$. So $\widehat W$ is appropriate only as a stand-in when $u_i$ is
unobserved, and even then report the signal-only version next to it. In the synthetic simulations,
where you \emph{know} $u_i$, use the true $W$. Never label the vote-total proxy simply ``welfare.''
\end{watchout}

\begin{definition}[Geographic-equity measures]
\label{def:equity-metrics}
For each rule report: the district funding shares
$f_d=\sum_{p\in g^M\cap\Pset^d}\cost(p)\big/\sum_{p\in g^M}\cost(p)$; the population shares
$\beta_d=n_d/n$; the smallest realised $\beta$ in the sense of Definition~\ref{def:geo},
$\min_d \tfrac{m}{\budget}\sum_{p\in g^M\cap\Pset^d}\cost(p)$ over districts with a valued project;
and the Gini coefficient of $(f_d)$. The theory predicts MES dominates on the realised $\beta$ and
on the gap $|f_d-\beta_d|$, with the DEM intermediate and the UAM worst on low-salience districts.
\end{definition}

% ─────────────────────────────────────────────────────────────────────
\subsection{Testable implications of the theory}
% ─────────────────────────────────────────────────────────────────────

Each diagnostic proposition implies something one can actually look for --- in generated data now, in
observed ballots later. Stating these turns the theory into hypotheses, which is what a referee will
want to see.

\begin{itemize}[nosep]
  \item \textbf{From Proposition~\ref{prop:noise} (noise).} Within a district, vote shares should track
        observable quality only weakly wherever noise dominates. Concretely: regress a project's
        within-district vote share $V(p)/\sum_{q\in\Pset^d}V(q)$ on cost, scope, and type,
        \emph{separately for each district}; a low $R^2$ in districts with many projects (large $k_d$)
        is the signature of noise domination.
  \item \textbf{From Proposition~\ref{prop:mob} (mobilization).} A district with a major institutional
        anchor should record more votes than its residents could cast alone:
        $V(\Pset^d)/(n_d\,\votes)>1$ for anchor districts.
  \item \textbf{From Proposition~\ref{prop:equity} (equity).} Funded budget per district should
        correlate \emph{negatively} with district need under the DEM (more marginal districts get less),
        and that correlation should move toward zero under the Method of Equal Shares. Choose the need
        proxy deliberately --- it is the same modelling choice as the equity benchmark of
        Section~\ref{sec:distortions}. Since we are \emph{not} using a CONAPO marginalization index, a
        natural stand-in is (inverse) salience, or a population-adjusted funding share; whatever the
        choice, state it and keep it consistent with the benchmark in Definition~\ref{def:geo}.
\end{itemize}

\begin{aside}\small
\textbf{Generating synthetic ballots.} For the simulation, draw each voter's district (or, under the
Method of Equal Shares, their approval scope) and within-scope weights from distributions calibrated to
plausible --- later, observed --- vote shares. A convenient generator is a Dirichlet-multinomial centred
on target within-district shares, with the approval ballot read off as $A_i=\{p:x_{i,p}\ge1\}$. When
real DEM vote totals (e.g.\ from a past \emph{presupuesto participativo}) become available, calibrate
the same generator to them; the synthetic exercise then becomes a genuine counterfactual --- run the
Method of Equal Shares on the calibrated ballots and compare the funded set, the welfare measures of
Definition~\ref{def:welfare}, and the equity measures of Definition~\ref{def:equity-metrics} against
what the DEM actually produced.
\end{aside}

% ─────────────────────────────────────────────────────────────────────
\subsection{Design levers and robustness}
% ─────────────────────────────────────────────────────────────────────

Vary, and report sensitivity to: the salience profile $(\salience_d)$ and how $\sigma_i$ is drawn;
the share of mobile (cross-district) voters; the intensity $c$ of institutional mobilization; the
spread of project costs relative to the per-voter share $\budget/n$ (this governs how much MES funds
directly versus in completion); and the completion rule. The qualitative ranking from
Section~\ref{sec:comparison} should be stable across these; any instability is itself a finding worth
reporting rather than hiding.

% ─────────────────────────────────────────────────────────────────────
\subsection{An implementation design for Durango}
% ─────────────────────────────────────────────────────────────────────

A concrete, defensible design to show the municipality:
\begin{enumerate}[label=\textbf{\arabic*.},nosep]
  \item \textbf{Ballot.} List all projects, grouped by district, with descriptions and costs; let
        voters approve any number --- no minimum, no maximum. Paper and digital both work.
  \item \textbf{Shares.} Split the budget equally among registered voters; each share is virtual,
        never real money.
  \item \textbf{Run.} Execute the Method of Equal Shares transparently on the collected ballots
        (open-source implementations exist, e.g.\ \texttt{equalshares.net}), and publish the
        round-by-round trace so the result is auditable.
  \item \textbf{Completion.} Spend any leftover greedily by approvals --- familiar, and it disposes of
        the residue the main rule cannot place. Flag that completion carries no proportionality
        guarantee, and design project cost tiers so the proportional phase reaches as far as
        possible.
\end{enumerate}

\begin{connection}
The pitch to Durango is not ``your system is broken'' but ``here is a transparent, auditable rule
that funds the same obvious winners while guaranteeing every neighbourhood a fair share by
construction --- and here is a simulation on plausible numbers showing the difference.'' That framing
is what a data-sharing conversation can be built on.
\end{connection}

% ═════════════════════════════════════════════════════════════════════
\section*{Reference shelf}
\addcontentsline{toc}{section}{Reference shelf}
% ═════════════════════════════════════════════════════════════════════

\noindent\textbf{Cite, do not reprove:}
\begin{itemize}[nosep]
  \item Myerson (1981); Mas-Colell, Whinston \& Green, \emph{Microeconomic Theory}, ch.~23 --- the
        Revelation Principle and Gibbard--Satterthwaite.
  \item Gibbard (1973); Satterthwaite (1975) --- the impossibility theorem itself.
  \item Peters, Pierczy\'nski \& Skowron (2021), ``Proportional Participatory Budgeting with
        Additive Utilities'' --- the Method of Equal Shares, EJR, and the EJR/strategy-proofness
        impossibility (their Thm.~4.8). The project site \texttt{equalshares.net} has
        implementations and intuition.
  \item Goel et al.\ (2019), ``Knapsack Voting'' --- the baseline PB model we extend.
  \item Martinelli (2006), ``Would Rational Voters Acquire Costly Information?'' --- background for the
        information function $\sigma$ and the noise result.
\end{itemize}

\bigskip
\noindent\textbf{Names recap.} \emph{Ours (coined here):} District-Exclusivity Mechanism (DEM);
Unrestricted Allocation Mechanism (UAM); district salience $\salience_d$; mobilizing institution and
vote displacement $\Delta_\mathcal{C}$; $\beta$-geographic equity; the signal/noise framing.
\emph{From the literature (cited):} Method of Equal Shares; Extended Justified Representation;
cohesive group; Gibbard--Satterthwaite; the Revelation Principle; DSIC/BIC.

\bigskip
\noindent\emph{The full corrected statements, proofs, and the review trail that motivated each
correction live in} \texttt{math-review/} \emph{and} \texttt{math-review/REVIEW-findings.md}.

\end{document}
